% =====================================================================
% Section 08, The four structural gaps
% =====================================================================

\section{Phase 1 \textbullet\ The Four Structural Gaps}
\label{sec:phase1-gaps}

The twelve research questions converge on \kw{four structural gaps}.
Each one is defined operationally (one formula), counted on the cleaned
data, and visualised in a specific Q-chart. They are the bridge between
Phase 1 (audit) and Phase 2 (verdict on new infrastructure): every Phase
2 hypothesis is an attempt to test whether the gap got smaller.

\subsection{G1 \textbullet\ Ghost Terminals · 115}

\textbf{Definition.}
\fn{Underused = (Total\_Routes > 0) AND (Boarding\_Near\_Terminal == 0)}
where \fn{Boarding\_Near\_Terminal} sums daily boardings within a
\stat{100 m} buffer of the terminal centroid.

\textbf{Counts.}
\begin{itemize}
 \item \stat{115} ghost terminals at the 100-m baseline.
 \item \stat{69} \kw{near-miss} (100--500 m), fix by relocating the stop.
 \item \stat{17} \kw{stop-too-far} (500 m--1 km), fix by improving
 walkability.
 \item \stat{29} \kw{truly isolated} ($>$1 km), audit and decommission.
\end{itemize}

\textbf{Recovery.} \stat{60\%} of ghosts disappear at a 500-m radius
(see Q11 recovery curve). The remaining 40\% are the genuine spatial
problem.

\textbf{Implication.} The ghost-terminal phenomenon is partly a
walkability problem (fixable) and partly a wasted-infrastructure
problem. The 29 truly isolated terminals are running zero passengers
yet still consuming routes that could be redirected.

\subsection{G2 \textbullet\ Empty-Return Terminals · 19}

\textbf{Definition.}
$\text{Empty Return Index} = 1 - \dfrac{\text{passengers}}{\text{vehicles}}$ within the terminal's 50-m flow buffer.
Critical threshold: \stat{$\geq 0.60$}.

\textbf{Counts.} \stat{19} terminals above the threshold. \stat{208} are healthy
(index = 0.0, every vehicle carries at least one passenger). The
distribution is  binary, see Q10.

\textbf{Magnitude.} The 19 critical terminals run \stat{193 empty
vehicle-trips per hour} between them. Zero passengers. Every day.

\textbf{Implication.} G2 is the operational-discipline gap. Routes are
being dispatched without demand on the return leg. Software-side
scheduling could close most of it; physical-infrastructure changes
cannot.

\subsection{G3 \textbullet\ Vehicle-Route Mismatch · 75\%}

\textbf{Definition.} Routes longer than 50 km whose assigned vehicle is
microbus (14 seats). Captures the systematic operational waste of
running a small vehicle on a long route.

\textbf{Counts.}
\begin{itemize}
 \item \stat{75\%} of routes $>$50 km use microbus.
 \item \stat{44} specific long-route violations identified.
 \item Vehicle mix: 861 microbus, 503 bus, 287 minibus, 115 tomnaya, 18 box.
\end{itemize}

\textbf{Cost manifestation.} The fare/km comparison from Q12 is the
direct reading: microbus = \stat{0.62 LE/km}, tomnaya = \stat{1.30 LE/km},
formal bus / metro = \stat{0.22--0.25 LE/km}. Riders pay the smaller
vehicle's price premium for trips that should have been served by a bus.

\textbf{Implication.} G3 is the procurement-and-allocation gap. Cost
minimisation drives the dispatch (microbuses are cheaper to run) at
the expense of service design.

\subsection{G4 \textbullet\ Underserved Hexagons · 79}

\textbf{Definition.}
$\text{Underserved\_Score} = \text{pct\_rank}(\text{Population}_{2018}) - \text{pct\_rank}(\text{Boarding total})$, clipped to $[0, 1]$. Threshold for ``underserved'': \stat{$> 0.5$}.

\textbf{Counts.} \stat{79 hexagons} cross the 0.5 threshold. \stat{381}
hexagons sit in the top quartile of the score (used by Phase 2 for
market sizing).

\textbf{Geography.} Imbaba, Shubra, Matariyya, the dense
informal-housing belt on the west and north of central Cairo. Layla
lives in this cluster.

\textbf{Implication.} G4 is the supply gap that everything else points
to. It is the demand the formal network does not visit.

\subsection{The Phase-1 \texttimes\ Phase-2 gap-closure matrix}

The four gaps are the rows of the synthesis matrix that Phase 2 will
fill with measured closure percentages (per-mode, within a 2-km
walkshed). Anticipating Section~\ref{sec:phase2-questions} (Q18b):

\rowcolors{2}{darkpanel}{darkbg}
\begin{table}[H]
\centering\small
\renewcommand{\arraystretch}{1.4}
\begin{tabular}{@{} >{\color{plotorange}\bfseries}l
 >{\color{plotyellow}\ttfamily}r
 >{\color{plotyellow}\ttfamily}r
 >{\color{plotyellow}\ttfamily}r @{}}
\toprule
\color{plotblue}Gap & \color{plotblue}Metro L3 & \color{plotblue}LRT & \color{plotblue}BRT \\
\midrule
G1 \ Ghosts & 15.7\% & 2.6\% & 7.8\% \\
G2 \ Empty-return & 5.3\% & 5.3\% & 10.5\% \\
G3 \ Vehicle mismatch & 0.0\% & 25.0\% & \stat{25.0\%} \\
G4 \ Underserved & 6.0\% & 1.6\% & 13.1\% \\
\bottomrule
\end{tabular}
\caption{Q18b \textbullet\ \% of Phase-1 gap sites within 2 km of a
post-2014 station. \stat{No cell above 25\%}, most below 16\%. The
single best cell (BRT $\times$ G3) reaches 25\%. The Cairo Monorail
would be a fourth column but it is not yet operational.}
\label{tab:gap-matrix}
\end{table}
\rowcolors{0}{}{}

\begin{callout}
\textbf{\color{plotblue}One sentence summary.}\quad The four gaps from
Phase 1 are real, measurable, and concentrated in specific districts;
the post-2014 infrastructure built since closes \emph{a quarter or
less} of any single gap. That single observation is what motivates
every hypothesis in Phase 2 and the product story in Section 16.
\end{callout}
